% Reusable preamble/header for Chinese research PDFs.
% Compile with XeLaTeX and keep ruc-logo.pdf beside the document.
\usepackage[a4paper,top=25mm,left=25mm,bottom=25mm,right=25mm,headheight=35pt,headsep=5mm,footskip=10mm]{geometry}
\usepackage{fontspec}
\usepackage{xeCJK}
\usepackage{graphicx}
\usepackage{fancyhdr}
\usepackage{eso-pic}
\usepackage{booktabs}
\usepackage{longtable}
\usepackage{array}
\usepackage{tabularx}
\usepackage{etoolbox}
\usepackage{enumitem}
\usepackage{ragged2e}
\usepackage{xcolor}
\usepackage{microtype}
\usepackage{needspace}
\usepackage{indentfirst}
\setmainfont{Times New Roman}
\setsansfont{Arial}
\setmonofont{Courier New}
\setCJKmainfont[AutoFakeBold=1.5]{Songti SC}
\setCJKsansfont[AutoFakeBold=1.5]{Songti SC}
\setCJKmonofont{Songti SC}
\newCJKfontfamily\songti[AutoFakeBold=1.5]{Songti SC}
\newCJKfontfamily\heiti[AutoFakeBold=2]{Songti SC}
\newCJKfontfamily\fangsong{Songti SC}
\setlength{\parindent}{2em}
\setlength{\parskip}{0pt}
\linespread{1.5}
\AtBeginDocument{\zihao{-4}\songti}
\setlength{\emergencystretch}{3em}
% Compact headings adapted from the user's reference paper; manual numbering remains in the titles.
\ctexset{
  section={format=\centering\songti\zihao{-3},beforeskip=1.0ex,afterskip=0.5ex},
  subsection={format=\raggedright\songti\zihao{-4},beforeskip=1.0ex,afterskip=0.5ex},
  subsubsection={format=\raggedright\songti\zihao{-4},beforeskip=1.0ex,afterskip=0.5ex}
}
\setlist{itemsep=0pt,topsep=0pt,parsep=0pt,partopsep=0pt}
\setlist[enumerate,1]{label=\arabic*.,leftmargin=2em}
\pagestyle{fancy}
\fancyhf{}
\fancyhead[C]{\includegraphics[pagebox=mediabox,width=5.2cm]{ruc-logo.pdf}}
\fancyfoot[C]{\zihao{-5}\thepage}
\renewcommand{\headrulewidth}{0.4pt}
\fancypagestyle{plain}{\fancyhf{}\fancyhead[C]{\includegraphics[pagebox=mediabox,width=5.2cm]{ruc-logo.pdf}}\fancyfoot[C]{\zihao{-5}\thepage}\renewcommand{\headrulewidth}{0.4pt}}
\fancypagestyle{empty}{\fancyhf{}\fancyhead[C]{\includegraphics[pagebox=mediabox,width=5.2cm]{ruc-logo.pdf}}\fancyfoot[C]{\zihao{-5}\thepage}\renewcommand{\headrulewidth}{0.4pt}}
\setlength{\tabcolsep}{5pt}
\renewcommand{\arraystretch}{1.0}
\AtBeginEnvironment{longtable}{\zihao{5}\linespread{1.0}\selectfont\setlength{\LTpre}{0.5\baselineskip}\setlength{\LTpost}{0.5\baselineskip}\renewcommand{\arraystretch}{1.0}}
\definecolor{linknavy}{RGB}{36,70,105}
\hypersetup{colorlinks=true,linkcolor=linknavy,urlcolor=linknavy,citecolor=linknavy}
\newcommand{\RUCTableCaption}[1]{\Needspace{3\baselineskip}\par\vspace{0.5\baselineskip}\begin{center}{\heiti\zihao{5}\textbf{#1}\par}\end{center}\vspace{0.5\baselineskip}}
\newcommand{\RUCTableAfterSpace}{\par\vspace{0.7\baselineskip}}
\newcommand{\RUCNote}[1]{\par\vspace{0.35\baselineskip}\noindent{\zihao{5}\songti\textbf{#1}\par}\vspace{0.7\baselineskip}}
% 参考文献条目正文与文章正文统一为小四号宋体；行距沿用全局1.5倍设置。
% 首条参考文献前增加半个正文行距，使标题与正文间距约为原来的两倍；
% 后续条目不增加额外间距。
\newif\ifRUCFirstReference
\RUCFirstReferencetrue
\newcommand{\RUCReference}[2]{\Needspace{8\baselineskip}\ifRUCFirstReference\vspace{0.5\baselineskip}\RUCFirstReferencefalse\fi\par\noindent\hangindent=2em\hangafter=1{{\zihao{-4}\songti #1\ #2}\par}}
\newcommand{\RUCSubtitle}{}
\newcommand{\RUCDateLine}{}
\makeatletter
\renewcommand{\maketitle}{%
  \par\begingroup\centering
  {\fontsize{22}{26.4}\selectfont\songti\textbf{\@title}\par}
  \ifx\RUCSubtitle\empty\else\vspace{10pt}{\fontsize{15}{18}\selectfont\songti\RUCSubtitle\par}\fi
  \ifx\RUCDateLine\empty\else\vspace{8pt}{\zihao{5}\songti\RUCDateLine\par}\fi
  \endgroup
  \thispagestyle{plain}
  \vspace{2pt}
}
\makeatother
